\begingroup
\makeatletter\tagsleft@true\let\veqno\@@leqno\makeatother
\clearpage\SourceLeaf{327b}\ArithmeticHead{320}{L}{ANALYSIS.\quad NACHLASS.}
\begin{center}\small\bfseries BEMERKUNGEN ZUR THEORIE DER TRANSZENDENTEN FUNKTIONEN.\end{center}
\par\indent Die hier als Hauptüberschrift gewählten Worte gehören in der Handschrift zu der den Abschnitt [I.] bildenden Abhandlung; wir haben diese mit den Aufzeichnungen zusammengefaßt, die in den Abschnitten [II.] und [III.] wiedergegeben sind, weil alle drei Stücke sich auf die Lehre von den elliptischen Funktionen beziehen.
\begin{center}\bfseries\textls[120]{Erläuterungen zu [I] und [II.], S. 287---311.}\end{center}
\par\indent Der art. [1.] enthält in seinem ersten Teile die Formel für die lineare Transformation der Thetafunktionen; es ist nämlich in neuerer Bezeichnungsweise (vergl. oben S. 275)
\[T=\sum_{k=-\infty}^{+\infty}e^{-a(k+\omega)^2}=e^{-a\omega^2}\vartheta_{00}\left(\frac{a\omega}{\pi i}\middle|e^{-a}\right)\]
\noindent und die Formel [1] lautet
\begin{equation}
e^{-a\omega^2}\vartheta_{00}\left(\frac{a\omega}{\pi i}\middle|e^{-a}\right)=\sqrt{\frac\pi a}\,\vartheta_{00}\left(\omega\middle|e^{-\frac{\pi^2}a}\right),\tag*{[1]$'$}
\end{equation}
\noindent also, wenn man $e^{-\frac{\pi^2}a}=q=e^{\tau\pi i}$ setzt (vergl. etwa \textsc{Weber}, \textit{Elliptische Funktionen und algebraische Zahlen}, Braunschweig 1891, S. 78, Gl. (11)),
\clearpage\SourceLeaf{328}\ArithmeticHead{321}{R}{{\scriptsize BEMERKUNGEN ZUR THEORIE DER TRANSZENDENTEN FUNKTIONEN.}}
\begin{equation}
e^{-\frac{\pi i\omega^2}\tau}\vartheta_{00}\left(\frac\omega\tau\middle|e^{-\frac{\pi i}\tau}\right)=\sqrt{-i\tau}\,\vartheta_{00}(\omega\mid e^{\tau\pi i}).\tag*{[1]$''$}
\end{equation}
\noindent $P$ bedeutet in den Gleichungen des art. [1.], wie auch sonst oft bei \textsc{Gauss} (siehe z. B. \textit{Disquisitiones arithmeticae}, Werke I, S. 413), den Umfang des Kreises vom Halbmesser Eins. Zu den Formeln selbst vergl. oben S. 223 die Formeln für $p\left(\frac1t,\frac{iu}t\right)$ und für $\sum h^{\alpha n^2+2\beta n+\gamma}$.
\par\indent Die den art. [1.] abschließende Bemerkung über das agM. enthält die auch im art. 16. der \textit{Determinatio attractionis} (1818), Werke III, S. 352 unten, zu findende \textsc{Landen}sche Transformation. Ihr Auftreten an dieser Stelle zeigt, daß \textsc{Gauss} die Bedeutung der Gleichung [1] als Transformationsformel wohl erkannt hat.
\par\indent Die artt. [2.], [3] geben die Umformung der Potenzreihen, deren Exponenten eine arithmetische Reihe zweiter Ordnung bilden (Thetanullreihen), in Produkte, vergl. schon in der Scheda Ac, oben S. 201, art. [9.]. Im art. [4.] wendet sich \textsc{Gauss} nach einer Zusammenfassung dieser Umformungsgleichungen zu den von den \textls[120]{beiden} Veränderlichen $x$ und $y$ abhängenden Reihen und Produkten (Thetafunktionen). Die Gleichung $\#$ ist jene schon in der Scheda Ac (siehe oben S. 204) auftretende berühmte Identität zwischen der Reihen- und der Produktdarstellung der Thetafunktion, vergl. auch die Abhandlung \textit{Zur Theorie der neuen Transscendenten}, Werke III, S. 446, Gl. 6. in Verbindung mit Gl. 9. auf S. 447 (diese Abhandlung stammt aus dem Handbuch 18, Bd, S. 221---223 und ist, da sie der Tagebuchaufzeichnung Nr. 139 entspricht, im Juni 1809 verfaßt). In den artt. [5.]---[10.] werden Transformationsformeln für die allgemeine Thetafunktion durch Induktion aufgestellt und dann in den artt. [11.], [12.] mit Hilfe des jetzt sogenannten \textsc{Hermite}schen Transformationsprinzips bewiesen, vergl. oben S. 277. Man vergl. zu diesen Formeln auch die bereits erwähnte Abhandlung \textit{Zur Theorie der neuen Transscendenten}, Werke III, S. 446 und die \textit{Hundert Theoreme}, ebenda S. 461ff. In der letzteren Abhandlung (die wohl aus der Zeit um 1825 stammt) werden für die drei geraden Thetafunktionen »besondere Functionalzeichen eingeführt« (Werke III, S. 465); es ist für $x=q$, $y=e^{2\pi vi}$
\[P(x,y)=\vartheta_{00}(v\mathord{\raisebox{.75ex}{\rule{.3pt}{.65ex}}}q),\quad Q(x,y)=\vartheta_{01}(v\mid q),\quad R(x,y)=\vartheta_{10}(v\mid q);\]
\noindent zu diesen tritt später (1827, Werke III, S. 472) das Zeichen $S(x,y)$ für die ungerade Thetafunktion. Wir fügen noch für einige der wichtigsten Formeln Verweisungen auf die entsprechenden Stellen der oben genannten späteren Abhandlungen hinzu.
\par\medskip
{\small\setlength{\tabcolsep}{2pt}\renewcommand{\arraystretch}{1.0}
\noindent\begin{tabular}{@{}lclcccp{6.4cm}@{}}
Im art. [3.],&zu&[2]&vergl.&Werke&III,&S. 447, Gl. 14. (auch Werke II, S. 20),\\
art. [4.],&»&[3]&»&»&»,&S. 447, Gl. 10.,\\
&»&[4]&»&»&»,&S. 447, Gl. 11.,\\
&»&[5]&»&»&»,&S. 446, Gl. 5.,\\
art. [7.],&»&[6]&»&»&»,&S. 449, Gl. 24.,\\
art. [8.],&»&[7]&»&»&»,&S. 450, Gl. 36.,\\
&»&[8]&»&»&»,&S. 451, Gl. 41.,\\
art. [9.],&»&1.&»&»&»,&S. 450, Gl. 29., für $-x$ statt $x$,\\
&»&2.&»&»&»,&S. 450, Gl. 30.,\quad »\qquad »\\
&»&[9]&»&»&»,&S. 449, Gl. 28.,\quad »\qquad »\\
&»&[10]&»&»&»,&S. 449, Gl. 25. und 27.
\end{tabular}\par}
\par\indent Im art. [13.] bedeutet $\varpi$ die lemniskatische Periode; zu dem daselbst gefundenen Maximalwert $\frac1{\sqrt[4]{2}}\sqrt{\frac\varpi\pi}$ wird man die aus der Scheda Aa stammenden numerischen Werte Werke III, S. 418 art. [8.] vergleichen können. In der Fußnote oben S. 302 ist $M=\log_{10}e$.
\par\medskip\noindent{\small X~1.\hfill41}\par
\clearpage\SourceLeaf{329}\ArithmeticHead{322}{L}{ANALYSIS.\quad NACHLASS.}
\par\indent Im art. [14.] kommt für die Transformationsformeln mit fünften Einheitswurzeln wieder das sog. \textsc{Hermite}sche Prinzip zur Anwendung; man wird diese Rechnungen als Vorbereitungen für die Fünfteilung anzusehen haben, die erst in der Abhandlung \textit{Zur Theorie der neuen Transscendenten}, Werke III, S. 456---459 vollständig ausgeführt wird. Die in der Fußnote am Schluß von art. [14.] befindliche Formel hat \textsc{Gauss} mit Bleistift an den Rand geschrieben, sie stimmt mit der Werke III, S. 475 (vom 29. August 1827, aus demselben Handbuch 16, Bb, S. 140 stammenden) für $5PP-pp$ überein.
\par\indent Siebenteilung und Dreiteilung werden in den artt. [15.], [16.] ausgeführt.
\par\indent Im Handbuch 16, Bb, folgen auf S. 51---53 jetzt noch Betrachtungen, die bereits Werke III, S. 442ff., art. [8.], [9.] ohne wesentliche Auslassungen abgedruckt sind und darum hier nicht noch einmal wiedergegeben werden; dagegen mußten des Zusammenhangs wegen, die art. [1.]---[7.] der Werke III, S. 436 beginnenden Abhandlung, die nur einen Auszug aus dem Text des Handbuchs 16, Bb darstellen, hier mit dem unverkürzten Texte aufs neue abgedruckt werden*).
\par\indent Die Abfassungszeit dieses Abschnitts [I.] ist, wie \textsc{Schering} (Werke III, S. 494, Absatz 3) andeutet, durch eine Mitteilung von \textsc{Gauss} an \textsc{Schumacher} auf das Jahr 1808 festgelegt. Mit dieser Mitteilung hat es folgende Bewandtnis. Zu der Zeit, wo \textsc{Jacobis} erster Brief an \textsc{Schumacher} (vom 13. Juni 1827, siehe \textsc{Jacobis} Werke I, S. 31) in Altona ankam, war \textsc{Gauss} gerade bei \textsc{Schumacher} zu Besuch (vergl. \textit{Wilhelm Olbers, Sein Leben und seine Werke} II, 2, Berlin 1909, S. 484, Fußnote). Während der Dauer von \textsc{Gauss}’ Anwesenheit und mit seiner »Approbation« schrieb \textsc{Schumacher} in seinem Antwortschreiben an \textsc{Jacobi}: »\textsc{Gauss} hat schon im Jahre 1808 die 3\,theilung, 5\,theilung und 7\,theilung entwickelt, und dabei die neuen darauf sich beziehenden Modulscalen gefunden« (siehe den \textit{Briefwechsel zwischen Gauss und Schumacher} II, S. 175, vergl. auch den ersten Brief \textsc{Jacobis} an \textsc{Legendre} vom 5. August 1827, \textsc{Jacobis} Werke I, S. 394, ferner \textsc{Legendre} ebenda S. 398, \textsc{Jacobi} ebenda S. 416, \textsc{Legendre} ebenda S. 418, 428). Nun finden wir in den art. [14.], [15.], [16.] unserer Abhandlung Untersuchungen über Fünfteilung, Siebenteilung und Dreiteilung, ferner hat \textsc{Gauss} am 6. August 1827, also bald nach seiner um den 18 Juli erfolgten Heimkehr von der Reise, auf S. 132 desselben Handbuchs 16, Bb, Aufzeichnungen begonnen, die an die hier in Rede stehenden Untersuchungen anknüpfen, dann allerdings sehr wesentlich darüber hinausgehn. Diese Aufzeichnungen, die sich bis zur S. 145 des Handbuchs erstrecken, sind Werke III, S. 470---480 abgedruckt. Die Angaben von \textsc{Gauss}, die \textsc{Schumacher} an \textsc{Jacobi} mitgeteilt hat, beziehen sich also sicher auf unsere Abhandlung [I.]. Weiter wird man annehmen können, daß der äußere Anstoß, der \textsc{Gauss} im Jahre 1808 zur Wiederaufnahme der Arbeiten an der Lehre von den elliptischen Funktionen geführt hat, durch den Brief \textsc{Schumachers} vom 2. April 1808 (siehe oben S. 242, [4.]) gegeben war. Vergl. den Abschnitt VI. des Aufsatzes »Über \textsc{Gauss}’ Arbeiten zur Funktionentheorie«.
\par\indent Auf den Seiten 63, 72, 73 des Handbuchs 16, Bb sind Formeln über lemniskatische Funktionen aufgezeichnet, die Werke III, S. 405, 406 untermischt mit älteren, und S. 409 Zeile 5 v. u. bis S. 412 untermischt mit späteren Aufzeichnungen wiedergegeben sind. Sie stehen, wie die Stelle Werke III, S. 412 zeigt, mit zahlentheoretischen Untersuchungen in Verbindung.
\par\indent Den Seiten 111---112 eben dieses Handbuchs ist der Abschnitt [II.] entnommen, auf den wir schon
\par\Needspace{12\baselineskip}\bigskip\noindent\rule{.25\linewidth}{.4pt}\par\smallskip
{\small\par\indent *) Eine von \textsc{Hattendorff} angefertigte \textls[120]{vollständige} Abschrift der Abhandlung [I.] \linebreak mit vielen dazu gehörigen Rechnungen und noch einigen andern auf den \textsc{Gauss}schen Nachlaß bezüglichen Aufzeichnungen ist nach \textsc{Hattendorffs} 1882 erfolgtem Ableben für das \textsc{Gauss}archiv erworben worden. Wir haben diese Abschrift beim Abdruck verglichen und konnten ihr einige Verbesserungen der Handschrift entnehmen. Vergl. K. \textsc{Hattendorff}, \textit{Die elliptischen Functionen in dem Nachlasse von Gauss}, Hannover, 1869. Eine Ableitung für die Werke III, S. 436---415 abgedruckten Formeln der hier in Rede stehenden \textsc{Gauss}schen Abhandlung gibt P. \textsc{Pepin} auf S. 94ff. seiner oben S. 282 genannten Schrift.\par}
\clearpage\SourceLeaf{330a}\ArithmeticHead{323}{R}{{\scriptsize BEMERKUNGEN ZUR THEORIE DER TRANSZENDENTEN FUNKTIONEN.}}
\noindent oben S. 279 hingewiesen haben*). Dieselbe Rechnung findet sich in etwas gedrängter Darstellung in den erwähnten am 6. August 1827 begonnenen Aufzeichnungen desselben Handbuchs, abgedruckt Werke III, S. 473, art. [5.]. Unser Abschnitt [II.] steht in dem Handbuche nach flächentheoretischen Betrachtungen, die Werke VIII, S. 405---407 abgedruckt sind und die, wie P. \textsc{Stäckel} a. a. O., S. 407 bemerkt, 1825 geschrieben sein dürften; in dieses Jahr wäre also auch der Abschnitt [II.] zu setzen. Daß \textsc{Gauss} aber diese Verifikationsrechnung schon 1800---1801 ausgeführt hat, zeigt nicht nur die Tagebuchnotiz Nr. 108 (vergl. die Erläuterungen zu dieser Notiz in dem unten folgenden Abdruck des \textit{Tagebuchs}), sondern auch der Umstand, daß der hier zur Anwendung gelangende Satz
\[ab'-ba'=\tfrac12ab(a^4-b^4),\]
\noindent von dem \textsc{Gauss} (oben S. 310) sagt, daß der Beweis davon tiefer liegt, mit dem \textit{Theorema} der Scheda Af (1801), oben S. 212, identisch ist**).
\par\bigskip\noindent\rule{.25\linewidth}{.4pt}\par\smallskip
{\small\par\indent *) Ein Auszug ist Werke III, S. 401, 402 abgedruckt.
\par\indent **) Auf S. 53 des Handbuchs 16, Bb bei der Werke III, S. 445 abgedruckten, auf die Differentialgleichungen für die Reihen $p,q$ bezüglichen Rechnung, steht dieser Satz in der Form
\[p^4-q^4=-4\frac{x\,dq}{q\,dx}+\frac{4x\,dp}{p\,dx}\]
\noindent mit einem Fragezeichen. Diese Formel fehlt in dem Abdruck Werke III, S. 445.\par}
\clearpage
\endgroup
